% GENERATED FILE. Edit canonical lessons, then run npm run generate:books.
% canonical-source-hash: fnv1a64:2822ce814de17e48
% canonical-lessons: PA-C289-dhund, PA-C289-kakkar, PA-C289-pala, PA-C289-tharna, PA-C289-sial

\chapter{The Cold Season Ahead}
\label{ch:pa-a2-the-cold-season-ahead}
\hlchaptermodality{289}

\begin{chapteropening}
\textbf{By the end of this chapter:} \emph{I can say what the cold months will be like.}

Complete the last lesson of chapter 289: I can say what the cold months will be like.
\end{chapteropening}

\section[\texorpdfstring{\pa{ਧੁੰਦ} (dhund)}{ਧੁੰਦ (dhund)}]{\pa{ਧੁੰਦ} (dhund) — fog; mist}

\label{lesson:PA-C289-dhund}



\begin{quote}
Before the new one: say the Punjabi for the hot season, then the Punjabi for a drought.
\end{quote}

\subsection*{You'll want to know: \pa{ਧੁੰਦ}}
\textbf{\pa{ਧੁੰਦ}} — \emph{dhund} — \textquotedblleft{}fog; mist\textquotedblright{}.

\begin{quote}
\pa{ਪੋਹ} \pa{ਵਿੱਚ} \pa{ਸਵੇਰੇ} \pa{ਧੁੰਦ} \pa{ਹੋਵੇਗੀ।}
\end{quote}

\begin{quote}
\emph{poh vich savere dhund hovegī.}
\end{quote}

\begin{quote}
In Poh there will be fog in the morning.
\end{quote}

A feminine noun.

\begin{cousinweb}
\textbf{\pa{ਧੁੰਦ}} (\emph{dhund}) is inherited from a reconstructed \emph{*dhūmāndha}, \textquotedblleft{}dim-sighted\textquotedblright{}; CDIAL lists Punjabi \emph{dhund}, \textquotedblleft{}mist, haziness\textquotedblright{}, under it (CDIAL 6858).
\end{cousinweb}

\begin{grammarlens}[title={What you've built}]
One more word for saying what comes next.
\end{grammarlens}

\subsection*{Your turn}
\begin{itemize}
\raggedright
  \item \emph{Say it:} \emph{dhund}
  \item \emph{Say it:} \emph{dhund}, once more
  \item \emph{Say it:} \emph{poh vich savere dhund hovegī.} — In Poh there will be fog in the morning.
  \item \emph{Recall:} say the Punjabi for the hot season, then the Punjabi for a drought, then say \emph{dhund} again
\end{itemize}

\begin{tcolorbox}[breakable,colback=teal!4,colframe=teal!35!black,title={Before you move on}]
What does \pa{ਧੁੰਦ} mean? (\textquotedblleft{}Fog; mist\textquotedblright{}.) Say it once more.
\end{tcolorbox}
\section[\texorpdfstring{\pa{ਕੱਕਰ} (kakkar)}{ਕੱਕਰ (kakkar)}]{\pa{ਕੱਕਰ} (kakkar) — frost}

\label{lesson:PA-C289-kakkar}



\begin{quote}
Before the new one: say the Punjabi for a drought, then the Punjabi for fog.
\end{quote}

\subsection*{You'll want to know: \pa{ਕੱਕਰ}}
\textbf{\pa{ਕੱਕਰ}} — \emph{kakkar} — \textquotedblleft{}frost\textquotedblright{}.

\begin{quote}
\pa{ਕੱਲ੍ਹ} \pa{ਸਵੇਰੇ} \pa{ਘਾਹ} \pa{ਉੱਤੇ} \pa{ਕੱਕਰ} \pa{ਹੋਵੇਗਾ।}
\end{quote}

\begin{quote}
\emph{kall savere ghāh utte kakkar hovegā.}
\end{quote}

\begin{quote}
Tomorrow morning there will be frost on the grass.
\end{quote}

A masculine noun.

\begin{cousinweb}
\textbf{\pa{ਕੱਕਰ}} (\emph{kakkar}) is inherited from Sanskrit \emph{karkara}, \textquotedblleft{}hard, firm\textquotedblright{}; CDIAL lists Punjabi \emph{kakkar}, \textquotedblleft{}frost\textquotedblright{}, under it (CDIAL 2819).
\end{cousinweb}

\begin{grammarlens}[title={What you've built}]
One more word for saying what comes next.
\end{grammarlens}

\subsection*{Your turn}
\begin{itemize}
\raggedright
  \item \emph{Say it:} \emph{kakkar}
  \item \emph{Say it:} \emph{kakkar}, once more
  \item \emph{Say it:} \emph{kall savere ghāh utte kakkar hovegā.} — Tomorrow morning there will be frost on the grass.
  \item \emph{Recall:} say the Punjabi for a drought, then the Punjabi for fog, then say \emph{kakkar} again
  \item \emph{Write it:} copy \textbf{\pa{ਕੱਕਰ}} once, checking the whole word
\end{itemize}

\begin{tcolorbox}[breakable,colback=teal!4,colframe=teal!35!black,title={Before you move on}]
What does \pa{ਕੱਕਰ} mean? (\textquotedblleft{}Frost\textquotedblright{}.) Say it once more.
\end{tcolorbox}
\section[\texorpdfstring{\pa{ਪਾਲਾ} (pālā)}{ਪਾਲਾ (pālā)}]{\pa{ਪਾਲਾ} (pālā) — bitter cold; a hard chill}

\label{lesson:PA-C289-pala}



\begin{quote}
Before the new one: say the Punjabi for fog, then the Punjabi for frost.
\end{quote}

\subsection*{You'll want to know: \pa{ਪਾਲਾ}}
\textbf{\pa{ਪਾਲਾ}} — \emph{pālā} — \textquotedblleft{}bitter cold; a hard chill\textquotedblright{}.

\begin{quote}
\pa{ਮਾਘ} \pa{ਵਿੱਚ} \pa{ਬਹੁਤ} \pa{ਪਾਲਾ} \pa{ਪੈਂਦਾ} \pa{ਹੈ।}
\end{quote}

\begin{quote}
\emph{māgh vich bahut pālā paindā hai.}
\end{quote}

\begin{quote}
In Magh it gets bitterly cold.
\end{quote}

A masculine noun.

\begin{cousinweb}
\textbf{\pa{ਪਾਲਾ}} (\emph{pālā}) is inherited from Sanskrit \emph{prālēya}, \textquotedblleft{}hail, frost, snow\textquotedblright{}; CDIAL lists Punjabi \emph{pālā}, \textquotedblleft{}frost, winter, cold, a chill\textquotedblright{}, under it (CDIAL 8959).
\end{cousinweb}

\begin{grammarlens}[title={What you've built}]
One more word for saying what comes next.
\end{grammarlens}

\subsection*{Your turn}
\begin{itemize}
\raggedright
  \item \emph{Say it:} \emph{pālā}
  \item \emph{Say it:} \emph{pālā}, once more
  \item \emph{Say it:} \emph{māgh vich bahut pālā paindā hai.} — In Magh it gets bitterly cold.
  \item \emph{Recall:} say the Punjabi for fog, then the Punjabi for frost, then say \emph{pālā} again
  \item \emph{Write it:} copy \textbf{\pa{ਪਾਲਾ}} once, checking the whole word
\end{itemize}

\begin{tcolorbox}[breakable,colback=teal!4,colframe=teal!35!black,title={Before you move on}]
What does \pa{ਪਾਲਾ} mean? (\textquotedblleft{}Bitter cold; a hard chill\textquotedblright{}.) Say it once more.
\end{tcolorbox}
\section[\texorpdfstring{\pa{ਠਰਨਾ} (ṭharnā)}{ਠਰਨਾ (ṭharnā)}]{\pa{ਠਰਨਾ} (ṭharnā) — to be chilled; to go numb with cold}

\label{lesson:PA-C289-tharna}



\begin{quote}
Before the new one: say the Punjabi for frost, then the Punjabi for bitter cold.
\end{quote}

\subsection*{You'll want to know: \pa{ਠਰਨਾ}}
\textbf{\pa{ਠਰਨਾ}} — \emph{ṭharnā} — \textquotedblleft{}to be chilled; to go numb with cold\textquotedblright{}.

\begin{quote}
\pa{ਪੋਹ} \pa{ਵਿੱਚ} \pa{ਹੱਥ} \pa{ਠਰਦੇ} \pa{ਹਨ।}
\end{quote}

\begin{quote}
\emph{poh vich hatth ṭharde han.}
\end{quote}

\begin{quote}
In Poh your hands go numb with cold.
\end{quote}

Past (I …): \pa{ਮੈਂ} \pa{ਬਹੁਤ} \pa{ਠਰਿਆ}, \emph{maiṁ bahut ṭhariā}.

\begin{cousinweb}
\textbf{\pa{ਠਰਨਾ}} (\emph{ṭharnā}) is inherited from a reconstructed \emph{*stharati}, \textquotedblleft{}is firm\textquotedblright{}; CDIAL lists Punjabi \emph{ṭharnā}, \textquotedblleft{}to be chilled or numb\textquotedblright{}, under it (CDIAL 13742).
\end{cousinweb}

\begin{grammarlens}[title={What you've built}]
One more word for saying what comes next.
\end{grammarlens}

\subsection*{Your turn}
\begin{itemize}
\raggedright
  \item \emph{Say it:} \emph{ṭharnā}
  \item \emph{Say it:} \emph{ṭharnā}, once more
  \item \emph{Say it:} \emph{poh vich hatth ṭharde han.} — In Poh your hands go numb with cold.
  \item \emph{Recall:} say the Punjabi for frost, then the Punjabi for bitter cold, then say \emph{ṭharnā} again
  \item \emph{Write it:} copy \textbf{\pa{ਠਰਨਾ}} once, checking the whole word
\end{itemize}

\begin{tcolorbox}[breakable,colback=teal!4,colframe=teal!35!black,title={Before you move on}]
What does \pa{ਠਰਨਾ} mean? (\textquotedblleft{}To be chilled; to go numb with cold\textquotedblright{}.) Say it once more.
\end{tcolorbox}
\section[\texorpdfstring{\pa{ਸਿਆਲ} (siāl)}{ਸਿਆਲ (siāl)}]{\pa{ਸਿਆਲ} (siāl) — winter; the cold season}

\label{lesson:PA-C289-sial}



\begin{quote}
Before the new one: say the Punjabi for bitter cold, then the Punjabi for to be chilled.
\end{quote}

\subsection*{You'll want to know: \pa{ਸਿਆਲ}}
\textbf{\pa{ਸਿਆਲ}} — \emph{siāl} — \textquotedblleft{}winter; the cold season\textquotedblright{}.

\begin{quote}
\pa{ਸਿਆਲ} \pa{ਵਿੱਚ} \pa{ਮੈਂ} \pa{ਕੰਬਲ} \pa{ਵਰਤਾਂਗਾ।}
\end{quote}

\begin{quote}
\emph{siāl vich maiṁ kambal vartāṅgā.}
\end{quote}

\begin{quote}
In winter I will use a blanket.
\end{quote}

A masculine noun.

\begin{cousinweb}
\textbf{\pa{ਸਿਆਲ}} (\emph{siāl}) is inherited from Sanskrit \emph{śītakāla}, \textquotedblleft{}cold season\textquotedblright{}; CDIAL lists Punjabi \emph{siāl}, \textquotedblleft{}winter\textquotedblright{}, under it (CDIAL 12486).
\end{cousinweb}

\begin{grammarlens}[title={What you've built}]
One more word for saying what comes next.
\end{grammarlens}

\subsection*{Your turn}
\begin{itemize}
\raggedright
  \item \emph{Say it:} \emph{siāl}
  \item \emph{Say it:} \emph{siāl}, once more
  \item \emph{Say it:} \emph{siāl vich maiṁ kambal vartāṅgā.} — In winter I will use a blanket.
  \item \emph{Recall:} say the Punjabi for bitter cold, then the Punjabi for to be chilled, then say \emph{siāl} again
  \item \emph{Write it:} copy \textbf{\pa{ਸਿਆਲ}} once, checking the whole word
\end{itemize}

\begin{tcolorbox}[breakable,colback=teal!4,colframe=teal!35!black,title={Before you move on}]
What does \pa{ਸਿਆਲ} mean? (\textquotedblleft{}Winter; the cold season\textquotedblright{}.) Say it once more.
\end{tcolorbox}
